\chapter{Phân tích và thiết kế hệ thống Yoca}
\label{Chapter3}

Chương này trình bày quá trình phân tích và thiết kế hệ thống Yoca, nền tảng web hỗ trợ phân tích dữ liệu on-chain, dữ liệu thị trường và hành vi ví blockchain. Tiếp nối bối cảnh và mục tiêu đã nêu ở Chương 1 cùng các hệ thống liên quan và cơ sở công nghệ đã trình bày ở Chương 2, Chương 3 tập trung làm rõ cách hệ thống Yoca được tổ chức để hiện thực hóa các mục tiêu đó, bao gồm ý tưởng ứng dụng, các tình huống sử dụng tiêu biểu, yêu cầu chức năng và phi chức năng, kiến trúc tổng thể, các luồng nghiệp vụ cốt lõi và định hướng thiết kế cơ sở dữ liệu. Do đề tài thuộc nhóm ứng dụng, chương chỉ tập trung vào các thành phần trọng tâm và các quyết định thiết kế quan trọng, không đi sâu mô tả toàn bộ chi tiết cài đặt hay toàn bộ bảng dữ liệu.

Yoca được xây dựng theo hướng ứng dụng web full-stack, trong đó frontend đảm nhiệm trình bày giao diện và trực quan hóa dữ liệu, còn backend đảm nhiệm xác thực, xử lý nghiệp vụ, tích hợp và chuẩn hóa dữ liệu từ bên thứ ba. Vì phải làm việc với dữ liệu blockchain có khối lượng lớn, biến động liên tục và phụ thuộc vào nhiều nhà cung cấp API, hệ thống cần có chiến lược truy xuất, lưu trữ đệm và đồng bộ dữ liệu hợp lý để đảm bảo khả năng phản hồi, ổn định và mở rộng.


\section{Ý tưởng ứng dụng và tình huống sử dụng}
\label{sec:y_tuong_ung_dung}

Yoca được phát triển với định hướng trở thành một nền tảng phân tích dữ liệu crypto và wallet intelligence, giúp người dùng theo dõi thị trường, phân tích token, quan sát hoạt động ví và khai thác các tín hiệu on-chain trong cùng một giao diện thống nhất. Trong bối cảnh thị trường blockchain phát triển nhanh, dữ liệu liên quan đến token, ví, giao dịch, dòng tiền và tin tức thường bị phân tán trên nhiều công cụ khác nhau. Người dùng muốn có một bức tranh đầy đủ thường phải tra cứu nhiều nguồn, đối chiếu dữ liệu thủ công và tự diễn giải các biểu đồ hoặc lịch sử giao dịch. Quá trình này tốn thời gian, dễ sai sót và gây khó khăn đối với người dùng chưa có nhiều kinh nghiệm phân tích dữ liệu on-chain.

Ý tưởng chính của Yoca là gom các lớp dữ liệu quan trọng vào một trải nghiệm phân tích thống nhất. Hệ thống không chỉ cung cấp thông tin thị trường như giá token, dữ liệu lịch sử và biểu đồ biến động, mà còn tập trung vào dữ liệu ở cấp độ ví như danh mục tài sản, lịch sử giao dịch, lịch sử chuyển token, lịch sử hoán đổi, biến động số dư và hiệu suất theo thời gian. Trên cơ sở đó, người dùng có thể quan sát một ví blockchain không chỉ như một địa chỉ lưu trữ tài sản, mà như một chuỗi hành vi tài chính có thể phân tích được.

Tình huống sử dụng đầu tiên là phân tích một ví blockchain. Người dùng nhập địa chỉ ví vào hệ thống để xem tổng quan danh mục, các token đang nắm giữ, giá trị ước tính, lịch sử giao dịch và các biểu đồ biến động theo thời gian. Thay vì đọc từng giao dịch riêng lẻ trên trình khám phá khối, người dùng có thể quan sát ví thông qua các góc nhìn tổng hợp như phân bổ tài sản, biến động số dư, số lượng giao dịch và các hoạt động swap đáng chú ý. Tình huống này đặc biệt hữu ích khi người dùng muốn hiểu một ví đang có xu hướng nắm giữ dài hạn, giao dịch thường xuyên hay thay đổi danh mục theo biến động thị trường.

Tình huống sử dụng thứ hai là theo dõi thị trường và token. Người dùng có thể tìm kiếm token, xem dữ liệu giá, dữ liệu lịch sử, các chỉ số thị trường và biểu đồ biến động. Chức năng này giúp người dùng theo dõi tài sản trong bối cảnh thị trường crypto có mức độ biến động cao. Khi dữ liệu token được kết hợp với dữ liệu ví, hệ thống có thể hỗ trợ người dùng hiểu rõ hơn sự thay đổi giá của một token ảnh hưởng như thế nào đến danh mục tài sản hoặc hiệu suất của ví.

Tình huống sử dụng thứ ba là so sánh và đánh giá nhiều ví. Đối với người dùng có nhu cầu theo dõi nhiều địa chỉ ví, hệ thống hỗ trợ quan sát và đối chiếu các chỉ số ở cấp độ ví. Người dùng có thể so sánh danh mục, mức độ hoạt động, hiệu suất và các dấu hiệu hành vi giữa các ví khác nhau. Đây là tình huống có giá trị đối với nhà nghiên cứu token, nhà phân tích ví hoặc người dùng muốn theo dõi các ví có ảnh hưởng trong hệ sinh thái.

Tình huống sử dụng thứ tư là thiết lập cảnh báo. Người dùng đã đăng nhập có thể cấu hình các cảnh báo liên quan đến token hoặc ví. Khi một điều kiện được thỏa mãn, chẳng hạn biến động giá, hoạt động ví hoặc một sự kiện đáng chú ý, hệ thống ghi nhận và hỗ trợ thông báo cho người dùng. Chức năng này làm giảm nhu cầu theo dõi thủ công liên tục, đồng thời giúp người dùng phản ứng nhanh hơn với các thay đổi quan trọng.

Tình huống sử dụng thứ năm là sử dụng AI để hỗ trợ diễn giải dữ liệu. Dữ liệu on-chain thường có nhiều trường kỹ thuật và không dễ đọc đối với người dùng phổ thông. Vì vậy, Yoca định hướng tích hợp trợ lý AI nhằm tóm tắt dữ liệu ví, giải thích các biến động nổi bật và hỗ trợ người dùng đặt câu hỏi trên dữ liệu đã được hệ thống chuẩn hóa. Trong thiết kế này, AI không thay thế các phép tính định lượng của hệ thống, mà đóng vai trò như một lớp hỗ trợ diễn giải kết quả bằng ngôn ngữ tự nhiên.


\section{Phân tích và đặc tả yêu cầu}
\label{sec:dac_ta_yeu_cau}

Dựa trên mục tiêu đề tài và các tình huống sử dụng đã nêu, hệ thống Yoca được phân tích theo hai nhóm yêu cầu chính: yêu cầu chức năng và yêu cầu phi chức năng. Yêu cầu chức năng mô tả những nghiệp vụ mà hệ thống cần hỗ trợ cho người dùng, trong khi yêu cầu phi chức năng mô tả các thuộc tính chất lượng như hiệu năng, bảo mật, tính ổn định, khả năng mở rộng và khả năng sử dụng.

\subsection{Các yêu cầu chức năng và phi chức năng}
% Liệt kê các chức năng chính: Token Analytics, Wallet Portfolio, Charting, AI Assistant. Yêu cầu phi chức năng: độ trễ < 0.5s.

Về yêu cầu chức năng, nhóm yêu cầu đầu tiên là quản lý tài khoản và xác thực người dùng. Hệ thống cần hỗ trợ đăng ký, đăng nhập, duy trì phiên làm việc và xác thực các thao tác yêu cầu quyền truy cập cá nhân hóa. Các chức năng như quản lý hồ sơ, watchlist, tag ví, cảnh báo và gói đăng ký chỉ nên được sử dụng bởi người dùng đã xác thực. Bên cạnh đăng nhập bằng tài khoản thông thường, hệ thống cũng hỗ trợ các hướng xác thực phù hợp với sản phẩm blockchain như đăng nhập thông qua ví.

Nhóm yêu cầu thứ hai là phân tích thị trường và token. Hệ thống cần cho phép người dùng tìm kiếm token, xem danh sách token, xem thông tin chi tiết token, dữ liệu giá, biểu đồ lịch sử và các chỉ số liên quan đến thị trường. Các dữ liệu này được lấy từ các nguồn bên thứ ba, sau đó được backend chuẩn hóa và trả về cho frontend dưới dạng phù hợp với bảng dữ liệu hoặc biểu đồ. Đây là nhóm chức năng tạo nền tảng để người dùng hiểu bối cảnh thị trường trước khi phân tích sâu hơn ở cấp độ ví.

Nhóm yêu cầu thứ ba là phân tích ví. Đây là một trong những nhóm chức năng trọng tâm của Yoca. Hệ thống cần cho phép người dùng nhập địa chỉ ví để xem tổng quan ví, danh mục token, lịch sử giao dịch, lịch sử chuyển tài sản, lịch sử swap, biến động số dư và hiệu suất theo thời gian. Đối với các ví có nhiều giao dịch, hệ thống cần có cơ chế phân trang, lưu trữ đệm và tái sử dụng dữ liệu để tránh việc phải truy xuất lại toàn bộ lịch sử từ nhà cung cấp bên ngoài trong mỗi lần người dùng mở trang.

Nhóm yêu cầu thứ tư là trực quan hóa dữ liệu. Hệ thống cần chuyển các dữ liệu đã xử lý thành biểu đồ dễ đọc, bao gồm biểu đồ giá token, biểu đồ giá trị danh mục, biểu đồ số dư, biểu đồ phân bổ tài sản, biểu đồ giao dịch và các biểu đồ hỗ trợ so sánh. Trực quan hóa không chỉ phục vụ mục đích trình bày mà còn là phương tiện chính giúp người dùng nhận diện xu hướng, biến động và các điểm bất thường trong dữ liệu.

Nhóm yêu cầu thứ năm là AI-assisted analysis và wallet chat. Hệ thống cần hỗ trợ chuyển dữ liệu ví hoặc dữ liệu token đã được xử lý thành ngữ cảnh có cấu trúc, sau đó sử dụng mô hình AI để tạo phần tóm tắt hoặc trả lời câu hỏi của người dùng. Kết quả AI cần được trình bày như một dạng hỗ trợ phân tích, không phải là kết luận tài chính tuyệt đối. Khi phù hợp, hệ thống có thể lưu lại kết quả phân tích để tránh lặp lại các tác vụ tốn chi phí và cải thiện tốc độ phản hồi trong các lần truy vấn sau.

Nhóm yêu cầu thứ sáu là cảnh báo. Người dùng đã xác thực cần có khả năng tạo, xem, cập nhật và xóa cảnh báo. Cảnh báo có thể liên quan đến biến động token, hoạt động ví hoặc các điều kiện do hệ thống hỗ trợ. Khi có dữ liệu mới phù hợp với điều kiện cảnh báo, hệ thống ghi nhận lịch sử và cung cấp thông tin cho người dùng. Đối với các cảnh báo liên quan đến hoạt động ví, hệ thống có thể sử dụng webhook hoặc cơ chế đồng bộ dữ liệu để phát hiện sự kiện mới.

Nhóm yêu cầu thứ bảy là quản lý hồ sơ, watchlist, nhãn ví và gói dịch vụ. Người dùng cần có khả năng lưu lại token hoặc ví quan tâm, gắn nhãn cho ví, quản lý thông tin cá nhân và theo dõi trạng thái gói đăng ký. Các chức năng này giúp hệ thống không chỉ là một công cụ tra cứu một lần, mà trở thành một không gian làm việc cá nhân hóa cho hoạt động phân tích dữ liệu blockchain.

Về yêu cầu phi chức năng, yêu cầu đầu tiên là hiệu năng. Do hệ thống phải làm việc với dữ liệu có khối lượng lớn và phụ thuộc vào nhiều API bên ngoài, backend cần ưu tiên đọc dữ liệu đã cache trước khi gọi lại nhà cung cấp. Đối với các dữ liệu có tần suất truy cập cao như dữ liệu token, dữ liệu biểu đồ hoặc dữ liệu lịch sử ví, cơ chế cache đóng vai trò quan trọng trong việc giảm độ trễ và giảm nguy cơ bị giới hạn tần suất truy vấn.

Yêu cầu phi chức năng thứ hai là tính ổn định. Hệ thống cần xử lý được các tình huống nhà cung cấp bên ngoài trả về lỗi, dữ liệu thiếu, bị giới hạn truy vấn hoặc có độ trễ cao. Khi gặp lỗi từ API bên ngoài, hệ thống cần có cách phản hồi rõ ràng cho frontend, đồng thời ưu tiên sử dụng dữ liệu đã lưu nếu dữ liệu đó vẫn còn phù hợp. Cách tiếp cận này giúp người dùng vẫn có thể tiếp tục sử dụng hệ thống trong nhiều trường hợp thay vì gặp lỗi hoàn toàn.

Yêu cầu phi chức năng thứ ba là bảo mật. Các chức năng liên quan đến tài khoản, watchlist, tag ví, cảnh báo, thanh toán và hồ sơ người dùng cần được bảo vệ bởi cơ chế xác thực và phân quyền. Token xác thực cần được quản lý an toàn, các API nhận dữ liệu từ người dùng cần được kiểm tra đầu vào, và các cấu hình nhạy cảm như API key, secret thanh toán hoặc secret xác thực không được đưa trực tiếp vào mã nguồn phía client.

Yêu cầu phi chức năng thứ tư là khả năng mở rộng. Mặc dù hệ thống hiện tập trung nhiều vào Solana, thiết kế cần đủ linh hoạt để có thể mở rộng sang các nguồn dữ liệu hoặc blockchain khác. Điều này đòi hỏi backend phải tách biệt logic nghiệp vụ với logic tích hợp từng provider. Khi cần bổ sung provider mới, hệ thống nên chỉ cần thêm adapter hoặc service tương ứng thay vì phải sửa trực tiếp nhiều phần giao diện và nghiệp vụ.

Yêu cầu phi chức năng thứ năm là khả năng sử dụng. Giao diện cần rõ ràng, nhất quán và phù hợp với người dùng có nhiều mức độ kinh nghiệm khác nhau. Người dùng mới cần có thể thực hiện các thao tác cơ bản như tìm token, nhập địa chỉ ví, đọc biểu đồ và xem danh mục mà không cần hiểu sâu về cấu trúc dữ liệu blockchain. Người dùng có kinh nghiệm hơn cần có đủ thông tin để phân tích sâu, so sánh dữ liệu và theo dõi các ví quan trọng.

\subsection{Sơ đồ Use Case tổng quát}

Sơ đồ use case tổng quát mô tả các tác nhân chính và nhóm chức năng lớn của hệ thống Yoca. Trong phạm vi thiết kế, hệ thống có bốn nhóm tác nhân chính: người dùng chưa đăng nhập, người dùng đã đăng nhập, quản trị viên và các dịch vụ bên thứ ba. Người dùng chưa đăng nhập có thể tra cứu một số thông tin công khai như thị trường, token hoặc ví. Người dùng đã đăng nhập có thể sử dụng các chức năng cá nhân hóa như quản lý watchlist, gắn nhãn ví, tạo cảnh báo, xem hồ sơ và sử dụng các chức năng yêu cầu xác thực. Quản trị viên có thể theo dõi, cấu hình hoặc quản lý một số dữ liệu hệ thống tùy theo phạm vi triển khai. Các dịch vụ bên thứ ba cung cấp dữ liệu blockchain, dữ liệu thị trường, AI, thanh toán và webhook.

% % TODO: Bổ sung hình sơ đồ Use Case tổng quát.
% \begin{figure}[htbp]
% \centering
% \fbox{
% \begin{minipage}{0.85\textwidth}
% \centering
% \vspace{1.2cm}
% Sơ đồ Use Case tổng quát của hệ thống Yoca \
% \textit{TODO: Chèn hình Use Case Diagram tại đây.}
% \vspace{1.2cm}
% \end{minipage}
% }
% \caption{Sơ đồ Use Case tổng quát của hệ thống Yoca}
% \label{fig}
% \end{figure}

% Hình~\ref{fig} cần thể hiện ranh giới hệ thống và các nhóm chức năng chính, bao gồm tra cứu thị trường, phân tích token, phân tích ví, so sánh ví, sử dụng AI Assistant, quản lý cảnh báo, quản lý hồ sơ và quản lý thanh toán hoặc gói đăng ký. Khi hoàn thiện sơ đồ, nhóm nên ưu tiên biểu diễn các nhóm use case chính thay vì liệt kê quá nhiều thao tác nhỏ, vì mục tiêu của báo cáo là làm rõ vai trò của từng tác nhân và các nghiệp vụ cốt lõi.

\section{Kiến trúc hệ thống tổng thể}
\label{sec:kien_truc_he_thong}

Yoca được tổ chức theo kiến trúc ứng dụng web full-stack, gồm hai ứng dụng chính là frontend và backend. Frontend được xây dựng bằng React và Vite, có nhiệm vụ hiển thị giao diện, nhận thao tác người dùng, gọi API và trực quan hóa dữ liệu. Backend được xây dựng bằng Hono, có nhiệm vụ định nghĩa API, xác thực yêu cầu, xử lý nghiệp vụ, truy cập cơ sở dữ liệu và tích hợp với các nhà cung cấp dữ liệu bên ngoài. Cơ sở dữ liệu PostgreSQL được sử dụng để lưu trữ dữ liệu người dùng, dữ liệu cache, metadata, lịch sử cảnh báo, dữ liệu thanh toán và các kết quả phân tích cần tái sử dụng.

Ở mức tổng thể, một yêu cầu từ người dùng thường đi qua các bước chính. Đầu tiên, người dùng thao tác trên một trang hoặc một component của frontend. Sau đó, frontend gọi một client service tương ứng để gửi yêu cầu đến backend. Backend tiếp nhận yêu cầu thông qua route, kiểm tra xác thực nếu cần, kiểm tra dữ liệu đầu vào và gọi domain service phù hợp. Domain service sẽ ưu tiên đọc dữ liệu trong PostgreSQL nếu dữ liệu đã được cache và còn phù hợp. Nếu dữ liệu chưa có hoặc đã quá cũ, service sẽ gọi nhà cung cấp bên ngoài, chuẩn hóa kết quả, lưu lại vào cơ sở dữ liệu và trả dữ liệu về cho frontend. Cuối cùng, frontend sử dụng dữ liệu này để hiển thị bảng, biểu đồ, báo cáo, cảnh báo hoặc phần trình bày từ AI.


% TODO: Bổ sung sơ đồ kiến trúc tổng thể.
\begin{figure}[H!]
    \centering
    \includegraphics[width=0.8\textwidth]{images/arch.png}
    \caption{Kiến trúc tổng thể của hệ thống Yoca}
    \label{fig}
\end{figure}

Hình~\ref{fig} nên thể hiện các thành phần chính gồm React frontend, Hono API backend, PostgreSQL, các domain service, các provider bên ngoài như Helius, Moralis, CoinGecko, Birdeye, Google Gemini, Stripe và Solana RPC. Ngoài ra, sơ đồ cũng nên làm rõ hướng dữ liệu đi từ frontend đến backend, từ backend đến cache, từ backend đến provider, sau đó quay lại frontend dưới dạng dữ liệu đã chuẩn hóa.



\subsection{Kiến trúc Monolith phân lớp và cơ chế Tải dữ liệu động (Lazy Loading)}

Trong phạm vi triển khai hiện tại, backend của Yoca được tổ chức theo hướng monolith phân lớp. Điều này có nghĩa là hệ thống backend được triển khai như một ứng dụng API thống nhất, nhưng bên trong được chia thành các lớp và miền nghiệp vụ riêng biệt. Các route chịu trách nhiệm tiếp nhận HTTP request và trả response. Các service chịu trách nhiệm xử lý nghiệp vụ. Các module tích hợp provider chịu trách nhiệm gọi API bên ngoài và chuẩn hóa dữ liệu. Lớp cơ sở dữ liệu chịu trách nhiệm định nghĩa schema, truy vấn và lưu trữ dữ liệu.

Cách tổ chức này phù hợp với quy mô của đồ án vì giúp hệ thống dễ triển khai, dễ kiểm thử và dễ theo dõi hơn so với việc tách thành nhiều microservices độc lập ngay từ đầu. Khi phát triển một sản phẩm có nhiều miền dữ liệu như token, ví, biểu đồ, AI, cảnh báo, thanh toán và hồ sơ người dùng, việc chia backend thành nhiều service vật lý có thể làm tăng độ phức tạp vận hành. Ngược lại, monolith phân lớp giúp nhóm tập trung vào việc hoàn thiện nghiệp vụ, đồng thời vẫn duy trì được cấu trúc module rõ ràng để có thể tách nhỏ trong tương lai nếu hệ thống phát triển lớn hơn.

Ở phía frontend, cơ chế tải dữ liệu động được áp dụng theo tương tác và ngữ cảnh sử dụng của người dùng. Thay vì tải toàn bộ dữ liệu ngay khi người dùng mở ứng dụng, mỗi trang hoặc thành phần giao diện chỉ gọi dữ liệu cần thiết cho chức năng hiện tại. Ví dụ, trang market chỉ cần dữ liệu thị trường và danh sách token, trang chi tiết token cần thêm dữ liệu lịch sử và biểu đồ, trong khi trang ví cần dữ liệu danh mục, giao dịch và lịch sử số dư. Cách tiếp cận này giúp giảm lượng dữ liệu tải ban đầu, cải thiện trải nghiệm người dùng và hạn chế số lượt gọi API không cần thiết.

Cơ chế tải dữ liệu động cũng phù hợp với đặc thù của dữ liệu blockchain. Một ví có thể có lịch sử giao dịch rất dài, nhưng người dùng không phải lúc nào cũng cần xem toàn bộ lịch sử đó ngay lập tức. Hệ thống có thể tải trước phần tổng quan, sau đó tải các phần chi tiết như giao dịch, swap hoặc biểu đồ số dư khi người dùng mở tab tương ứng. Cách làm này giúp phân bổ chi phí truy vấn theo nhu cầu thực tế, tránh gây tải lớn cho backend và provider bên ngoài.

\subsection{Chiến lược lưu trữ đệm (Caching) đa tầng}
% Rất quan trọng: Giải thích cơ chế TTL/window filtering cho dữ liệu ví, và on-conflict upserts cho dữ liệu thị trường (Token market/chart caches).

Caching là một trong những quyết định thiết kế quan trọng của Yoca. Dữ liệu on-chain và dữ liệu thị trường thường được lấy từ các nhà cung cấp bên thứ ba, trong khi các nhà cung cấp này có thể áp dụng giới hạn tần suất truy vấn, giới hạn gói sử dụng hoặc thay đổi độ trễ theo thời điểm. Nếu mỗi lần người dùng mở một trang hoặc thay đổi bộ lọc hệ thống đều gọi trực tiếp provider, ứng dụng sẽ dễ gặp lỗi rate limit, phản hồi chậm và khó kiểm soát chi phí. Vì vậy, hệ thống cần lưu trữ đệm dữ liệu đã lấy và tái sử dụng trong những trường hợp phù hợp.

Chiến lược caching của Yoca được thiết kế theo hướng ưu tiên cơ sở dữ liệu. Đối với các dữ liệu có thể tái sử dụng như lịch sử ví, danh mục tài sản, giao dịch, chuyển token, swap, dữ liệu biểu đồ và dữ liệu token, backend sẽ kiểm tra dữ liệu trong PostgreSQL trước. Nếu dữ liệu đã tồn tại và còn nằm trong ngưỡng hợp lệ, hệ thống trả về dữ liệu cache. Nếu dữ liệu chưa có hoặc không còn đủ mới, hệ thống gọi provider bên ngoài để đồng bộ phần dữ liệu cần thiết, sau đó chuẩn hóa và lưu lại vào cơ sở dữ liệu.

Đối với dữ liệu ví, caching cần quan tâm đến khái niệm khoảng thời gian dữ liệu đã được đồng bộ. Một ví có thể đã được đồng bộ lịch sử giao dịch cho một giai đoạn nhất định, nhưng chưa đầy đủ ở một giai đoạn khác. Vì vậy, hệ thống cần sử dụng metadata hoặc thông tin coverage để phân biệt phần dữ liệu đã có và phần dữ liệu còn thiếu. Cách tiếp cận này giúp hệ thống không phải tải lại toàn bộ lịch sử ví mỗi lần người dùng truy vấn, đồng thời vẫn đảm bảo khả năng bổ sung dữ liệu mới khi cần.

Đối với dữ liệu token và dữ liệu thị trường, hệ thống có thể sử dụng cơ chế TTL và ngưỡng cập nhật. Dữ liệu thị trường gần thời gian thực cần được cập nhật thường xuyên hơn, trong khi dữ liệu lịch sử cũ có thể được tái sử dụng lâu hơn. Khi nhận dữ liệu mới, backend có thể thực hiện upsert để cập nhật hoặc bổ sung bản ghi mà không tạo dữ liệu trùng lặp. Cách tiếp cận này giúp duy trì tập dữ liệu rolling ổn định cho các biểu đồ và bảng thống kê.

Caching không chỉ phục vụ hiệu năng mà còn giúp hệ thống ổn định hơn. Khi provider bên ngoài gặp lỗi tạm thời, dữ liệu cache có thể được sử dụng như một phương án dự phòng nếu vẫn còn phù hợp. Điều này đặc biệt quan trọng đối với các trang phân tích ví và biểu đồ, vì người dùng thường kỳ vọng hệ thống vẫn hiển thị được dữ liệu đã truy xuất trước đó thay vì báo lỗi hoàn toàn khi provider không phản hồi.



\subsection{Cơ chế đồng bộ hóa dữ liệu từ các nhà cung cấp (Helius, CoinGecko)}

Yoca tích hợp nhiều nhà cung cấp dữ liệu để phục vụ các miền nghiệp vụ khác nhau. Đối với dữ liệu ví Solana, hệ thống sử dụng Helius để truy xuất danh mục, lịch sử giao dịch, chuyển token và swap. Đối với dữ liệu thị trường và dữ liệu lịch sử giá token, hệ thống sử dụng CoinGecko. Ngoài ra, hệ thống còn có các đường tích hợp khác như Moralis cho một số luồng dữ liệu EVM hoặc dữ liệu dự phòng, Birdeye như một hướng fallback hoặc legacy path cho một số dữ liệu token, Google Gemini cho AI-assisted analysis, Stripe cho thanh toán và Solana RPC hoặc Helius RPC cho xác minh thanh toán on-chain.

Việc đồng bộ dữ liệu từ nhiều provider đòi hỏi backend phải có lớp chuẩn hóa dữ liệu. Mỗi provider có định dạng phản hồi, giới hạn truy vấn và cách biểu diễn dữ liệu khác nhau. Nếu frontend hoặc domain service phụ thuộc trực tiếp vào dữ liệu thô từ provider, hệ thống sẽ khó bảo trì và dễ bị ảnh hưởng khi provider thay đổi. Vì vậy, Yoca sử dụng các adapter hoặc service chuyên biệt để gọi provider, chuyển dữ liệu về DTO nội bộ và chỉ trả về định dạng đã được chuẩn hóa cho các lớp phía trên.

Đối với dữ liệu ví, luồng đồng bộ thường bắt đầu bằng việc người dùng nhập địa chỉ ví hoặc mở lại một ví đã từng phân tích. Backend kiểm tra metadata cache để xác định dữ liệu đã có. Nếu dữ liệu thiếu hoặc cần cập nhật, backend gọi provider để lấy dữ liệu mới theo từng loại như portfolio, transactions, transfers hoặc swaps. Sau đó, dữ liệu được chuẩn hóa, lưu vào các bảng cache tương ứng và trả về cho frontend. Cách tiếp cận này giúp hệ thống quản lý dữ liệu ví theo từng miền nhỏ thay vì xử lý toàn bộ dữ liệu ví như một khối đơn lẻ.

Đối với dữ liệu token, backend cần ánh xạ token address sang định danh phù hợp với nguồn dữ liệu giá. Đây là bước quan trọng vì mỗi provider có thể dùng khóa định danh khác nhau. Sau khi ánh xạ được token, hệ thống có thể lấy dữ liệu thị trường, dữ liệu giá lịch sử hoặc dữ liệu biểu đồ. Các dữ liệu này sau đó được lưu lại để phục vụ các lần truy vấn tiếp theo.

Cơ chế đồng bộ cũng cần xử lý lỗi và giới hạn truy vấn. Khi provider trả về lỗi 429 hoặc lỗi tạm thời, hệ thống cần có chính sách chờ, retry hoặc trả về dữ liệu cache nếu có. Khi nhiều API key được cấu hình, hệ thống có thể xoay vòng key để giảm nguy cơ gián đoạn. Những cơ chế này giúp hệ thống hoạt động ổn định hơn trong điều kiện phụ thuộc vào nhiều dịch vụ bên ngoài.



\section{Thiết kế chi tiết các luồng nghiệp vụ cốt lõi (Sequence Diagrams)}
\label{sec:thiet_ke_luong_nghiep_vu}

Phần này trình bày ba luồng nghiệp vụ cốt lõi của hệ thống: truy xuất và tổng hợp lịch sử giao dịch ví, truy vấn và trực quan hóa dữ liệu biểu đồ, và vấn đáp dữ liệu thông qua AI Assistant. Các luồng này đại diện cho những vấn đề kỹ thuật chính của Yoca, bao gồm truy xuất dữ liệu on-chain, sử dụng cache, chuẩn hóa dữ liệu, hiển thị biểu đồ và diễn giải dữ liệu.

\subsection{Luồng truy xuất và tổng hợp lịch sử giao dịch ví}

Luồng truy xuất và tổng hợp lịch sử giao dịch ví bắt đầu khi người dùng nhập một địa chỉ ví hoặc mở trang chi tiết ví. Frontend gửi yêu cầu đến backend thông qua service tương ứng. Backend tiếp nhận yêu cầu tại route ví, kiểm tra định dạng địa chỉ, xác định loại dữ liệu cần truy xuất và gọi wallet domain service.

Wallet service trước hết kiểm tra dữ liệu cache trong PostgreSQL. Nếu hệ thống đã có dữ liệu tổng quan, danh mục và lịch sử giao dịch còn phù hợp, service sẽ đọc dữ liệu từ cache và trả về cho route. Nếu cache chưa tồn tại hoặc chưa đủ dữ liệu cho khoảng thời gian người dùng yêu cầu, service gọi lớp fetcher để truy xuất dữ liệu từ provider bên ngoài. Dữ liệu nhận được từ provider được chuẩn hóa thành định dạng nội bộ, sau đó được lưu vào các bảng cache và metadata tương ứng.

Sau khi dữ liệu đã được chuẩn hóa, backend tổng hợp các thông tin cần thiết cho frontend. Các thông tin này có thể bao gồm danh sách token trong ví, giá trị ước tính, lịch sử giao dịch, giao dịch chuyển token, hoạt động swap và các thống kê tổng quan. Frontend nhận dữ liệu trả về, cập nhật trạng thái giao diện và hiển thị dữ liệu dưới dạng bảng hoặc biểu đồ.

% % TODO: Bổ sung sequence diagram cho luồng truy xuất ví.
% \begin{figure}[htbp]
% \centering
% \fbox{
% \begin{minipage}{0.85\textwidth}
% \centering
% \vspace{1.2cm}
% Sequence Diagram: Luồng truy xuất và tổng hợp lịch sử giao dịch ví \
% \textit{TODO: Chèn sơ đồ tuần tự tại đây.}
% \vspace{1.2cm}
% \end{minipage}
% }
% \caption{Luồng truy xuất và tổng hợp lịch sử giao dịch ví}
% \label{fig}
% \end{figure}

% Hình~\ref{fig} nên thể hiện các thành phần chính gồm User, React Frontend, Hono Route, Wallet Service, Cache Retriever, Provider Fetcher, PostgreSQL và Helius. Điểm cần nhấn mạnh trong sơ đồ là backend không gọi provider ngay lập tức trong mọi trường hợp, mà kiểm tra cache trước, sau đó chỉ gọi provider khi dữ liệu thiếu hoặc đã hết hạn.



\subsection{Luồng truy vấn và trực quan hóa biểu đồ (Chart Data Flow)}

Luồng truy vấn dữ liệu biểu đồ bắt đầu khi người dùng mở một trang có biểu đồ hoặc thay đổi tham số như token, ví, khoảng thời gian hoặc loại biểu đồ. Frontend gọi chart API thông qua service đã định nghĩa. Request gửi lên backend chứa các thông tin cần thiết để xác định loại biểu đồ, đối tượng phân tích và khoảng thời gian cần truy vấn.

Backend tiếp nhận request tại chart route. Tùy theo loại biểu đồ, route có thể gọi chart service, token service hoặc wallet service. Ví dụ, biểu đồ lịch sử giá token sẽ dựa nhiều vào dữ liệu token và dữ liệu thị trường, trong khi biểu đồ biến động số dư ví sẽ dựa vào dữ liệu ví và dữ liệu lịch sử số dư. Service tương ứng kiểm tra cache trước, sau đó gọi provider nếu dữ liệu chưa đủ hoặc cần cập nhật.

Dữ liệu biểu đồ sau khi được lấy từ cache hoặc provider sẽ được chuẩn hóa thành chuỗi thời gian. Mỗi điểm dữ liệu cần có mốc thời gian, giá trị và các metadata cần thiết để frontend có thể hiển thị. Đối với một số biểu đồ phức tạp, chẳng hạn biểu đồ số dư theo token hoặc biểu đồ giá trị danh mục, backend có thể cần kết hợp dữ liệu giao dịch với dữ liệu giá lịch sử để tính ra giá trị tại từng thời điểm.

Frontend nhận dữ liệu đã chuẩn hóa và truyền vào các component biểu đồ. Các component này chịu trách nhiệm hiển thị trục thời gian, đường dữ liệu, tooltip, chú thích, trạng thái tải dữ liệu và trạng thái lỗi. Cách tách biệt giữa backend xử lý dữ liệu và frontend hiển thị biểu đồ giúp hệ thống dễ mở rộng thêm loại biểu đồ mới mà không làm thay đổi quá nhiều cấu trúc tổng thể.

% % TODO: Bổ sung sequence diagram cho Chart Data Flow.
% \begin{figure}[htbp]
% \centering
% \fbox{
% \begin{minipage}{0.85\textwidth}
% \centering
% \vspace{1.2cm}
% Sequence Diagram: Luồng truy vấn và trực quan hóa biểu đồ \
% \textit{TODO: Chèn sơ đồ tuần tự tại đây.}
% \vspace{1.2cm}
% \end{minipage}
% }
% \caption{Luồng truy vấn và trực quan hóa dữ liệu biểu đồ}
% \label{fig}
% \end{figure}

% Hình~\ref{fig} nên thể hiện các thành phần gồm Chart Component, Chart API Client, Hono Chart Route, Chart Service, Token hoặc Wallet Service, PostgreSQL Cache và provider như CoinGecko hoặc Helius. Điểm cốt lõi của luồng là dữ liệu biểu đồ phải được chuẩn hóa trước khi trả về frontend để đảm bảo các component biểu đồ có thể tái sử dụng.

\subsection{Luồng vấn đáp dữ liệu thông qua AI Assistant}

Luồng vấn đáp thông qua AI Assistant bắt đầu khi người dùng đặt câu hỏi hoặc yêu cầu hệ thống phân tích một ví, một token hoặc một tập dữ liệu đã hiển thị. Frontend gửi câu hỏi cùng với ngữ cảnh cần thiết đến backend. Backend kiểm tra xác thực nếu chức năng yêu cầu người dùng đăng nhập, sau đó gọi service phụ trách wallet chat hoặc AI analysis.

Trước khi gọi mô hình AI, backend cần chuẩn bị ngữ cảnh dữ liệu. Đây là bước quan trọng vì AI không nên tự tạo ra số liệu hoặc tự suy đoán dữ liệu không có trong hệ thống. Service sẽ thu thập các chỉ số đã được tính toán, dữ liệu ví đã chuẩn hóa, dữ liệu token liên quan hoặc các kết quả phân tích có sẵn. Sau đó, backend chuyển các dữ liệu này thành một ngữ cảnh có cấu trúc để gửi đến mô hình AI.

Mô hình AI sử dụng ngữ cảnh đã chuẩn bị để tạo câu trả lời bằng ngôn ngữ tự nhiên. Câu trả lời có thể bao gồm tóm tắt danh mục ví, nhận xét về biến động tài sản, các giao dịch đáng chú ý hoặc giải thích một biểu đồ nhất định. Sau khi nhận kết quả, backend có thể lưu lại câu trả lời hoặc phiên chat nếu chức năng lưu lịch sử được bật. Frontend hiển thị câu trả lời trong giao diện chat hoặc trong vùng phân tích tương ứng.

% % TODO: Bổ sung sequence diagram cho AI Assistant.
% \begin{figure}[htbp]
% \centering
% \fbox{
% \begin{minipage}{0.85\textwidth}
% \centering
% \vspace{1.2cm}
% Sequence Diagram: Luồng vấn đáp dữ liệu thông qua AI Assistant \
% \textit{TODO: Chèn sơ đồ tuần tự tại đây.}
% \vspace{1.2cm}
% \end{minipage}
% }
% \caption{Luồng vấn đáp dữ liệu thông qua AI Assistant}
% \label{fig}
% \end{figure}

% Hình~\ref{fig} nên thể hiện các thành phần gồm User, Chat Component, Hono Route, Wallet Chat Service, PostgreSQL, AI Provider và giao diện trả kết quả. Điểm quan trọng cần thể hiện là AI chỉ nhận dữ liệu đã được backend chuẩn hóa và tóm tắt, thay vì truy cập trực tiếp vào dữ liệu thô hoặc tự tạo ra số liệu.


\section{Thiết kế Cơ sở dữ liệu}
\label{sec:thiet_ke_csdl}

Cơ sở dữ liệu của Yoca được thiết kế nhằm phục vụ ba mục tiêu chính. Mục tiêu thứ nhất là lưu trữ dữ liệu người dùng và các thông tin cá nhân hóa như tài khoản, ví liên kết, watchlist, nhãn ví, cảnh báo và gói đăng ký. Mục tiêu thứ hai là lưu trữ dữ liệu cache phục vụ phân tích ví, token, biểu đồ và kết quả AI. Mục tiêu thứ ba là lưu trữ metadata để quản lý độ mới, khoảng dữ liệu đã đồng bộ và trạng thái xử lý của các luồng dữ liệu phụ thuộc provider bên ngoài.

Thiết kế cơ sở dữ liệu của hệ thống không nên được trình bày bằng cách giải thích toàn bộ từng bảng, từng cột, vì điều đó dễ biến báo cáo thành tài liệu kỹ thuật chi tiết. Thay vào đó, phần này tập trung vào các nhóm bảng trọng tâm và vai trò của chúng trong kiến trúc tổng thể. Các chi tiết cụ thể như tên cột, kiểu dữ liệu và ràng buộc có thể được trình bày trong phụ lục nếu cần.

\subsection{Sơ đồ thực thể mối quan hệ (ERD)}

Sơ đồ ERD của Yoca cần thể hiện các nhóm thực thể chính thay vì liệt kê quá chi tiết mọi cột. Ở mức tổng quan, cơ sở dữ liệu có thể được chia thành các nhóm: nhóm người dùng và xác thực, nhóm ví và dữ liệu ví, nhóm token và dữ liệu thị trường, nhóm biểu đồ và cache, nhóm cảnh báo, nhóm AI/chat, nhóm thanh toán và nhóm tin tức hoặc ngữ cảnh token.

% % TODO: Bổ sung ERD thật từ Drizzle schema.
% \begin{figure}[htbp]
% \centering
% \fbox{
% \begin{minipage}{0.85\textwidth}
% \centering
% \vspace{1.2cm}
% Sơ đồ ERD tổng quan của hệ thống Yoca \
% \textit{TODO: Chèn ERD tại đây.}
% \vspace{1.2cm}
% \end{minipage}
% }
% \caption{Sơ đồ thực thể mối quan hệ tổng quan của hệ thống Yoca}
% \label{fig}
% \end{figure}

% Hình~\ref{fig} nên ưu tiên thể hiện mối quan hệ giữa các nhóm bảng trọng tâm. Nhóm người dùng liên kết với các bảng hồ sơ, ví liên kết, watchlist, nhãn ví, cảnh báo và gói đăng ký. Nhóm ví liên kết với các bảng cache về tổng quan ví, danh mục, giao dịch, chuyển token, swap, số dư và phân tích AI. Nhóm token liên kết với dữ liệu thị trường, dữ liệu lịch sử và dữ liệu biểu đồ. Nhóm cảnh báo liên kết với điều kiện cảnh báo, lịch sử gửi và trạng thái xử lý. Nhóm AI/chat liên kết với phiên chat, tin nhắn và kết quả phân tích đã lưu.

\subsection{Thiết kế chi tiết các bảng trọng tâm}
% Trình bày các bảng chính phục vụ caching như wallet_history, token_metadata...

Nhóm bảng người dùng và xác thực lưu trữ thông tin tài khoản, phương thức đăng nhập, ví liên kết và các dữ liệu cá nhân hóa. Đây là nhóm bảng làm nền cho các chức năng yêu cầu xác thực như quản lý watchlist, tag ví, cảnh báo, hồ sơ cá nhân và thanh toán. Việc tách dữ liệu xác thực khỏi dữ liệu phân tích giúp hệ thống dễ kiểm soát quyền truy cập và tránh trộn lẫn dữ liệu cá nhân với dữ liệu blockchain công khai.

Nhóm bảng ví là nhóm bảng quan trọng nhất đối với chức năng wallet intelligence. Nhóm này lưu các dữ liệu như tổng quan ví, danh mục token, giao dịch, chuyển token, swap, lịch sử số dư, PnL hoặc kết quả phân tích. Một số bảng trong nhóm này đóng vai trò cache dữ liệu lấy từ provider, trong khi một số bảng khác đóng vai trò metadata để xác định dữ liệu đã được đồng bộ đến đâu. Nhờ đó, hệ thống có thể tái sử dụng dữ liệu đã có và chỉ truy xuất thêm phần còn thiếu khi người dùng yêu cầu.

Nhóm bảng token và thị trường lưu trữ dữ liệu liên quan đến token, ánh xạ token address với định danh bên ngoài, dữ liệu giá hiện tại, dữ liệu lịch sử và các thông tin hỗ trợ biểu đồ. Nhóm bảng này giúp hệ thống giảm phụ thuộc vào việc gọi API thị trường trong mỗi lần người dùng mở trang token hoặc biểu đồ. Đối với dữ liệu gần thời gian thực, hệ thống có thể áp dụng TTL để quyết định khi nào cần cập nhật lại.

Nhóm bảng biểu đồ và cache lưu trữ các chuỗi dữ liệu đã được tính toán hoặc chuẩn hóa để frontend có thể hiển thị nhanh hơn. Trong một số trường hợp, dữ liệu biểu đồ có thể được tính từ dữ liệu ví, dữ liệu token và dữ liệu giá lịch sử. Việc lưu lại kết quả trung gian giúp giảm chi phí xử lý lại, đặc biệt với các biểu đồ được truy cập nhiều lần hoặc yêu cầu tính toán trên khoảng thời gian dài.

Nhóm bảng cảnh báo lưu trữ định nghĩa cảnh báo, điều kiện kích hoạt, trạng thái, lịch sử xử lý và lịch sử gửi thông báo. Đối với các cảnh báo liên quan đến ví, hệ thống có thể kết hợp dữ liệu webhook hoặc dữ liệu đồng bộ theo chu kỳ để phát hiện sự kiện mới. Việc lưu lịch sử cảnh báo giúp người dùng xem lại các sự kiện đã xảy ra và giúp hệ thống tránh gửi trùng thông báo.

Nhóm bảng AI và chat lưu trữ phiên chat, tin nhắn, ngữ cảnh phân tích và kết quả phản hồi từ mô hình AI. Do các tác vụ AI có thể tốn chi phí và phụ thuộc vào provider bên ngoài, việc lưu lại kết quả trong một số trường hợp giúp hệ thống giảm số lần gọi lặp lại. Đồng thời, dữ liệu đã lưu cũng giúp người dùng tiếp tục phiên phân tích trước đó thay vì phải bắt đầu lại từ đầu.

Nhóm bảng thanh toán và gói đăng ký lưu trữ thông tin về gói dịch vụ, lịch sử thanh toán và trạng thái subscription. Hệ thống có thể hỗ trợ cả luồng thanh toán thông qua nhà cung cấp thanh toán truyền thống và luồng xác minh thanh toán on-chain. Vì các dữ liệu này liên quan đến quyền sử dụng chức năng và lịch sử tài chính của người dùng, chúng cần được xử lý cẩn trọng và chỉ truy cập thông qua các API đã xác thực.

Nhóm bảng tin tức và ngữ cảnh token lưu trữ các bài viết, nguồn tin, batch tin tức và các dữ liệu ngữ cảnh hỗ trợ phân tích token hoặc biến động thị trường. Dữ liệu tin tức không phải là thành phần cốt lõi như dữ liệu ví hoặc dữ liệu token, nhưng đóng vai trò bổ sung để người dùng hiểu bối cảnh bên ngoài của các biến động giá hoặc hoạt động thị trường.

Nhìn chung, thiết kế cơ sở dữ liệu của Yoca được định hướng theo các miền nghiệp vụ thay vì gom tất cả dữ liệu vào một số bảng lớn. Cách tổ chức này giúp hệ thống dễ mở rộng, dễ kiểm soát dữ liệu cache và phù hợp với cấu trúc service ở backend. Khi triển khai thực tế, các bảng chi tiết có thể thay đổi theo quá trình phát triển, nhưng các nhóm dữ liệu trọng tâm nêu trên vẫn là nền tảng chính cho kiến trúc lưu trữ của hệ thống.

\section{Kế hoạch kiểm thử và tiêu chí chấp nhận}
\subsection{Phạm vi và các mức kiểm thử}
\subsection{Dữ liệu và kịch bản kiểm thử}
\subsection{Tiêu chí chấp nhận hệ thống}

\section{Tổng kết chương}

Chương này đã trình bày quá trình phân tích và thiết kế hệ thống Yoca từ góc nhìn nghiệp vụ và kiến trúc. Hệ thống được định hướng như một nền tảng web hỗ trợ phân tích dữ liệu on-chain, dữ liệu thị trường và hành vi ví blockchain trong một giao diện thống nhất. Các tình huống sử dụng chính bao gồm phân tích ví, theo dõi token, so sánh ví, thiết lập cảnh báo và sử dụng AI để hỗ trợ diễn giải dữ liệu.

Về mặt yêu cầu, Yoca cần đáp ứng các nhóm chức năng liên quan đến xác thực, phân tích thị trường, phân tích ví, trực quan hóa dữ liệu, AI Assistant, cảnh báo, hồ sơ người dùng và thanh toán. Đồng thời, hệ thống cần đảm bảo các yêu cầu phi chức năng về hiệu năng, ổn định, bảo mật, khả năng mở rộng và khả năng sử dụng.

Về mặt kiến trúc, hệ thống được tổ chức theo mô hình ứng dụng web full-stack với React frontend, Hono backend và PostgreSQL làm cơ sở dữ liệu chính. Backend được thiết kế theo hướng monolith phân lớp, chia thành các miền nghiệp vụ rõ ràng như wallet, token, chart, auth, alert, AI, payment và profile. Cơ chế caching đóng vai trò quan trọng trong việc giảm phụ thuộc vào provider bên ngoài, cải thiện hiệu năng và tăng độ ổn định của hệ thống.

Các luồng nghiệp vụ cốt lõi như truy xuất lịch sử ví, truy vấn dữ liệu biểu đồ và vấn đáp thông qua AI Assistant đã được phân tích ở mức thiết kế. Đây là cơ sở để Chương 4 tiếp tục trình bày kết quả cài đặt, môi trường triển khai, các chức năng đã hoàn thiện và quá trình đánh giá hệ thống.